%=====================================================================
%  Appendices A-C
%
%  The published article numbers appendix tables and figures continuously
%  with the main text (Table 7 is followed by Table A.8; Fig. 8 by
%  Fig. B.9). The counters are therefore set explicitly below instead of
%  being reset at each appendix section.
%=====================================================================

% Cells of the literature-review tables are set ragged right, as in the
% published article.
\newcolumntype{R}[1]{>{\raggedright\arraybackslash}p{#1}}

% Appendix tables/figures carry the appendix letter as a prefix: "A.8", not
% "Appendix A.8" (\thesection spells out "Appendix A" in elsarticle).
\renewcommand{\thetable}{\Alph{section}.\arabic{table}}
\renewcommand{\thefigure}{\Alph{section}.\arabic{figure}}

%---------------------------------------------------------------------
\section{Literature review}\label{app:literature}
%---------------------------------------------------------------------
\setcounter{table}{7}   % next table becomes A.8

See Tables~\ref{tab:litreview1} and~\ref{tab:litreview2}.

{\scriptsize
\setlength{\tabcolsep}{3pt}
\renewcommand{\arraystretch}{1.05}
\begin{longtable}{@{}
    R{0.085\linewidth}
    R{0.075\linewidth}
    R{0.135\linewidth}
    R{0.095\linewidth}
    R{0.145\linewidth}
    R{0.095\linewidth}
    R{0.085\linewidth}
    R{0.185\linewidth}@{}}
  \caption{Summary of studies exploring cyclists' behavior and subjective safety
    at shared paths (Part 1).}\label{tab:litreview1}\\
  \toprule
  Study & Study area & Objective & Data collection method & Attributes
    & Behavior Indicators & Subjective Indicators & Findings\\
  \midrule
  \endfirsthead
  \multicolumn{8}{@{}l}{\textit{Table \thetable\ (continued)}}\\
  \toprule
  Study & Study area & Objective & Data collection method & Attributes
    & Behavior Indicators & Subjective Indicators & Findings\\
  \midrule
  \endhead
  \bottomrule
  \endlastfoot

  \multicolumn{8}{@{}l}{\textbf{Shared path}}\\
  \addlinespace[2pt]

  \citet{kazemzadeh2024forwhom}
    & Lund, Sweden
    & Explore cyclists' comfort during interactions in shared spaces,
      considering demographic heterogeneity
    & Intercept survey (594 cyclists)
    & 1. Cyclist demographics (age/gender/education)
      2. Cycling experience
      3. Interaction type (passing/meeting)
    &
    & Comfort
    & 1. Female cyclists, older adults, and less experienced cyclists generally
      report lower comfort in shared spaces
      2. Significant demographic differences affect comfort perception,
      highlighting the need for targeted infrastructure improvements.\\
  \addlinespace[4pt]

  \citet*{vonstulpnagel2024matter}
    & Berlin, Germany
    & Explore how environmental and demographic factors affect cyclists'
      perceived safety
    & Survey (21,500 participants)
    & 1. Buffer type between cycling and pedestrian lane
      (None/Greens/Cobbled)
      2. Sidewalk usage (Free/Cafes/shops)
    &
    & Perceived safety
    & Cyclists feel safe on sidewalks shared with pedestrians, at the expense of
      pedestrians' subjective safety\\
  \addlinespace[4pt]

  \citet*{berghoefer2023prefer}
    & Braun\-schweig (Campus), Germany
    & Explore factors affecting cyclists' infrastructure preferences and
      subjective safety
    & Cycling simulator (39 participants)
    & Advisory lane/Cycle path/Footpath:
      1. Pedestrian volume (5/25 [pedestrians])
      2. Vehicle volume (200/1000 [veh/h])
      Street/Multi-use path:
      1. Pedestrian volume (5/25 [pedestrians])
      2. Vehicle volume (100/400 [veh/h])
    &
    & 1. Mental Comfort
      2. Interaction
      3. Environment
      4. Ease of Use
      5. Physical Comfort
    & 1. Cycle paths that are separated from other road users were most
      preferred, while crowded footpaths and uphill segments were least
      preferred
      2. Cyclists riding on paths/lanes involving pedestrians felt less safe and
      comfortable when pedestrian volume was higher.\\
  \addlinespace[4pt]

  \citet{batista2023exploring}
    & Braun\-schweig, Germany
    & Explore factors affecting pedestrians' and cyclists' infrastructure
      preferences
    & Survey (408 participants)
    & 1. Space configuration (central placement of design
      elements/segregation types)
      2. Design elements (benches/fountains/bol\-lards/play elements)
      3. Presence of motor vehicles
    & 1. Crossing path choice
      2. Interaction behaviors
    & 1. Comfort
      2. Enjoyment level
      3. Preference for design attributes
    & 1. Both pedestrians and cyclists favor centrally located design elements
      and protective barriers.
      2. Motor vehicles are perceived negatively by both pedestrians and
      cyclists.
      3. Effective shared space designs incorporate integration, placemaking,
      and informal segregation.\\
  \addlinespace[4pt]

  \citet{zhang2023space}
    & Global (Systematic Review)
    & Explore interactions and safety between pedestrians and micro-mobility
      vehicles (MMVs) in shared spaces
    & Systematic literature review (87 publications)
    & 1. Infrastructure features
      2. Road user characteristics
      3. Operation condition
    & 1. Swerving
      2. Passing
      3. Overtaking
      4. Lateral displacement
    & 1. Perceived safety
      2. Comfort level
      3. Attitudes toward pedestrians
    & Key factors include sufficient infrastructure width, clear regulations,
      and speed management. Effective solutions include user education and
      infrastructure improvements.\\
  \addlinespace[4pt]

  \citet{pashkevich2022visual}
    & Kraków, Poland
    & Explore visual attention and speeds of pedestrians, cyclists, and electric
      scooter (ES) riders
    & Field experiment with eye tracker (12 participants)
    & 1. Visual attention distribution (roadway/roadside/all three road users)
      2. Speeds and interaction behaviors
    & 1. Fixation frequency and duration
      2. Interaction behaviors (maneuvering)
    & 1. Perceived safety
      2. Comfort level
    & 1. Riders (bicycles and ES) mainly focus visual attention on roadway and
      pedestrians, highlighting hazard awareness.
      2. Pedestrians tend to look more at surroundings.
      3. Cyclists and ES riders travel at similar speeds (16 km/h), suggesting
      ES should be regulated similarly to bicycles for safety.\\
  \addlinespace[4pt]

  \citet{che2021interaction}
    & Singapore
    & Explore the relationship among perception, attitude, and behavior of
      cyclists and pedestrians
    & 1. Survey for subjective indicators (300 cyclists)
      2. Video observation for behavior indicator (Before keeping left rule:
      119 cyclists After keeping left rule: 59 cyclists)
    & 1. On-coming pedestrian-cyclist passing
      2. Cyclist overtaking pedestrian
      3. Rule of keep left for all road users
    & Lane-keeping action:
      1. Keep left
      2. Keep right
      3. Zigzag
    & 1. Perceived safety
      2. Perceived control
      3. Perceived priority
    & 1. Cyclists' perceived control, perceived risk level, and perceived
      priority influence their attitude and further affect their behavior
      2. Cyclists tend to compensate for pedestrians' risks, thereby indicating
      a more cooperative way of space-sharing\\
  \addlinespace[4pt]

  \citet{huang2021investigating}
    & Hefei (Campus), China
    & Explore how behavior of cyclists and different movements of pedestrian
      crowds affect each other
    & Closed-field experiment (95 participants)
    & 1. Pedestrian and cyclist density (flow)
      2. Pedestrian and cyclist flow direction
      (homogeneous/unidirec\-tional/bidirectional)
    & 1. Speed
      2. Pedestrian lane formation
    &
    & 1. Cyclists travel slower when pedestrian density is higher
      2. Pedestrian lanes were more concentrated in bidirectional flow than
      unidirectional flow\\
  \addlinespace[4pt]

  \citet{gkekas2020perceived}
    & Vancouver (Campus), Canada
    & Explore factors affecting the risk of pedestrian-cyclist interactions
    & Survey (337 participants)
    &
    &
    & 1. Perceived safety
      2. Comfort
    & 1. Both pedestrians and cyclists considered crowding and pedestrian
      inattention as key factors contributing to incidents
      2. Cyclists feel safer when paths get wider\\
  \addlinespace[4pt]

  \citet{stulpnagel2020gaze}
    & Freiburg, Germany
    & Explore cyclists' gaze behavior relates to subjective risk perception and
      local spatial configurations
    & Field experiment with eye-tracking (20 participants)
    & 1. Spatial complexity
      2. Subjective risk perception
    & 1. Sight vector length
      2. Gaze angle
      3. Fixation duration
    & 1. Perceived safety
    & 1. Cyclists' gaze behavior is strongly influenced by spatial openness and
      complexity; open areas increase perceived risk.
      2. Experienced cyclists demonstrated longer sight vectors and narrower
      gaze angles, indicating more focused attention and hazard anticipation.
      3. Recommendations for urban planning include balancing visibility and
      spatial structure to manage visual workload and enhance cyclists'
      perceived safety.\\
\end{longtable}
}

{\scriptsize
\setlength{\tabcolsep}{3pt}
\renewcommand{\arraystretch}{1.05}
\begin{longtable}{@{}
    R{0.085\linewidth}
    R{0.075\linewidth}
    R{0.135\linewidth}
    R{0.095\linewidth}
    R{0.145\linewidth}
    R{0.095\linewidth}
    R{0.085\linewidth}
    R{0.185\linewidth}@{}}
  \caption{Summary of studies exploring cyclists' behavior and subjective safety
    at shared paths and unsignalized intersections (Part 2).}%
  \label{tab:litreview2}\\
  \toprule
  Study & Study area & Objective & Data collection method & Attributes
    & Behavior Indicators & Subjective Indicators & Findings\\
  \midrule
  \endfirsthead
  \multicolumn{8}{@{}l}{\textit{Table \thetable\ (continued)}}\\
  \toprule
  Study & Study area & Objective & Data collection method & Attributes
    & Behavior Indicators & Subjective Indicators & Findings\\
  \midrule
  \endhead
  \bottomrule
  \endlastfoot

  \citet*{nikiforiadis2019can}
    & Thessa\-loniki, Greece
    & Explore factors influencing pedestrian perception of comfort and safety on
      shared sidewalks and pedestrian street
    & 1. Survey (300 pedestrians)
      2. Field measurements
    & 1. Pedestrian and bicycle lane widths
      2. Pedestrian and bicycle volumes and speeds
    &
    & 1. Comfort
      2. Perceived safety
      3. Space sufficiency
    & 1. Pedestrian comfort and safety negatively impacted by increasing bicycle
      lane width relative to pedestrian area
      2. Pedestrian perceived level of service strongly influenced by space
      sufficiency and pedestrian unit flow rate, emphasizing the importance of
      appropriate spatial design.\\
  \addlinespace[4pt]

  \citet{beitel2018assessing}
    & Montreal (Campus), Canada
    & Explore factors affecting the risk of pedestrian-cyclist interactions
    & 1. Video observation for speed (2,200 cyclists)
      2. Survey for conflict rate (872 responses)
    & Pedestrian density (flow)
    & 1. Speed
      2. Conflict rate
      3. Angle of approach
      4. Post-encroachment time
    &
    & 1. Cyclists travel slower when pedestrian density is higher
      2. Conflict rate increase when pedestrian density is higher\\
  \addlinespace[4pt]

  \citet{boufous2018impact}
    & Sydney, Australia
    & Explore how environmental factors affect cycling speed
    & Video observation (5,421 cyclists)
    & 1. Width ($\leq$3.5/$>$3.5 m)
      2. Pedestrian flow ($<$20 /20-100/100-199/$\geq$200)
      3. Centerline (with/without)
      4. Commuter path (with/without)
      5. Visual segregation (yes/no)
    & Speed
    &
    & Cyclists travel slower on shared paths with smaller widths, higher
      pedestrian flow, centerline, and visual segregation\\
  \addlinespace[4pt]

  \citet{essa2018road}
    & Vancouver, Canada
    & Explore factors affecting the risk of pedestrian-cyclist interactions
    & Video observation (940 cyclists, 800 streets)
    & 1. Pedestrian density (flow)
      2. Dismount sign
    & 1. Speed
      2. Conflict rate
    &
    & 1. Cyclists travel slower when pedestrian density is higher
      2. Conflict rate decreases when bicycle dismount sign is presented\\
  \addlinespace[4pt]

  \multicolumn{8}{@{}l}{\textbf{Unsignalized intersection}}\\
  \addlinespace[2pt]

  \citet{mohammadi2024understanding}
    & Gothenburg, Sweden
    & Explore factors affecting cyclist behavior by interacting with vehicles
    & Cycling simulator (25 participants)
    & 1. The difference in time to arrival at the intersection (DTA)
      (1.2/2.5/3.5 [s])
      2. Intersection visibility (IV) (22/27 [m])
      3. Pedaling
      4. Looking duration
    & 1. Average speed
      2. Braking distance
      3. Yielding decision
    &
    & 1. Cyclist's DTA and IV significantly affect braking distance
      2. Cyclist's pedaling, looking duration, and DTA have a significant impact
      on cyclist's yielding behavior\\
  \addlinespace[4pt]

  \citet{mohammadi2023how}
    & Gothenburg, Sweden
    & Explore factors affecting cyclist behavior by interacting with vehicles
    & Video observation (105 interactions)
    & 1. Pedaling
      2. Head movement
      3. Hand gesture
      4. Initial speed
      5. The difference in time to arrival at the intersection (DTA)
    & Yielding decision
    &
    & Cyclist's pedaling, head movement, DTA, and speed of cyclist and vehicle
      have a significant impact on cyclist's yielding behavior\\
  \addlinespace[4pt]

  \citet{ng2017cyclist}
    & Queensland, Australia
    & Explore which types of cycling infrastructure cyclists perceive to be the
      safest at unsignalized intersections
    & Survey (214 Participants)
    & 12 different types of cycling infrastructure with each type having three
      different interactions with vehicles
    &
    & Perceived safety
    & 1. Off-road bicycle paths and footpaths were perceived to be the safest
      cycling infrastructure at unsignalized intersections
      2. Cyclists' perceived safety is influenced by motorists' yielding
      behavior toward them\\
\end{longtable}
}

%---------------------------------------------------------------------
\section{Perceived unsafety and stress}\label{app:unsafetystress}
%---------------------------------------------------------------------
\setcounter{figure}{8}   % next figure becomes B.9

See Figs.~\ref{fig:b9} and~\ref{fig:b10}.

\begin{figure}[htbp]
  \centering
  \includegraphics[width=0.86\linewidth]{figB9_unsafety_stress_shared.pdf}
  \caption{Perceived unsafety and stress in eight scenarios at the shared path
    environment.}
  \label{fig:b9}
\end{figure}

\begin{figure}[htbp]
  \centering
  \includegraphics[width=0.86\linewidth]{figB10_unsafety_stress_unsig.pdf}
  \caption{Perceived unsafety and stress in eight scenarios at the unsignalized
    intersection environment.}
  \label{fig:b10}
\end{figure}

%---------------------------------------------------------------------
\section{Model estimation results with all variables}\label{app:allvariables}
%---------------------------------------------------------------------
\setcounter{table}{9}   % next table becomes C.10

See Tables~\ref{tab:c10}--\ref{tab:c14}.

\begin{table}[htbp]
  \centering
  \caption{Model estimation results for participant average speed at the shared
    path.}
  \label{tab:c10}
  \small
  \begin{threeparttable}
  \begin{tabular}{@{}l r r r l@{}}
    \toprule
    Variables & Coefficient & Std. error & t-value & p-value\\
    \midrule
    \multicolumn{5}{c}{Participant average speed [km/hour]}\\
    \addlinespace[2pt]
    Intercept                                        & 17.001 & 1.206 & 14.094 & \pvl\ ***\\
    Path width high                                  &  3.146 & 0.563 &  5.584 & \pvl\ ***\\
    Pedestrian flow high                             & -1.907 & 0.563 & -3.385 & 0.001 **\\
    Bicycle flow high                                & -1.050 & 0.563 & -1.864 & 0.063\\
    Road width high $\times$ Pedestrian flow high    &  0.314 & 0.651 &  0.482 & 0.630\\
    Path width high $\times$ Bicycle flow high       & -0.941 & 0.651 & -1.446 & 0.149\\
    Pedestrian flow high $\times$ Bicycle flow high  &  1.870 & 0.651 &  2.874 & 0.004 **\\
    \midrule
    \multicolumn{2}{@{}l}{Adjusted R-squared} & \multicolumn{3}{l}{0.485}\\
    \bottomrule
  \end{tabular}
  \begin{tablenotes}[flushleft]\footnotesize
    \item \sigstars
  \end{tablenotes}
  \end{threeparttable}
\end{table}

\begin{table}[htbp]
  \centering
  \caption{Model estimation results for perceived unsafety at the shared path.}
  \label{tab:c11}
  \small
  \begin{threeparttable}
  \begin{tabular}{@{}l r r r l@{}}
    \toprule
    Variables & Odds ratio & Std. error & t-value & p-value\\
    \midrule
    \multicolumn{2}{@{}l}{Intercepts (Thresholds)}
      & \multicolumn{3}{l}{Perceived unsafety [Likert]}\\
    \addlinespace[2pt]
    Perceived unsafety 1 vs 2 (or higher)             &   0.018 & 1.119 & -3.575 & \pvl\ ***\\
    Perceived unsafety 2 (or lower) vs 3 (or higher)  &   0.355 & 1.097 & -0.944 & 0.345\\
    Perceived unsafety 3 (or lower) vs 4 (or higher)  &   3.033 & 1.101 &  1.088 & 0.313\\
    Perceived unsafety 4 (or lower) vs 5              & 175.550 & 1.235 &  4.185 & \pvl\ ***\\
    \addlinespace[2pt]
    Participant average speed                         &   0.892 & 0.041 & -2.798 & 0.005 **\\
    Road Width high                                   &   0.655 & 0.436 & -0.972 & 0.331\\
    Bicycle flow high                                 &   4.982 & 0.434 &  3.697 & \pvl\ ***\\
    Pedestrian flow high                              &   3.947 & 0.422 &  3.254 & 0.001 **\\
    Fixation on pedestrians (\%)                      &   0.996 & 0.014 & -0.274 & 0.784\\
    Fixation on bicycles (\%)                         &   1.003 & 0.030 &  0.099 & 0.921\\
    Road Width high $\times$ Pedestrian flow high     &   0.437 & 0.474 & -1.749 & 0.080\\
    Road Width high $\times$ Bicycle flow high        &   0.327 & 0.476 & -2.374 & 0.018 *\\
    Bicycle flow high $\times$ Pedestrian flow high   &   0.530 & 0.457 & -1.335 & 0.182\\
    \midrule
    \multicolumn{2}{@{}l}{Akaike information criterion (AIC)}
      & \multicolumn{3}{l}{682.75}\\
    \bottomrule
  \end{tabular}
  \begin{tablenotes}[flushleft]\footnotesize
    \item \sigstars
  \end{tablenotes}
  \end{threeparttable}
\end{table}

\begin{table}[htbp]
  \centering
  \caption{Model estimation results for participant average speed at the
    unsignalized intersection.}
  \label{tab:c12}
  \small
  \begin{threeparttable}
  \begin{tabular}{@{}l r r r l@{}}
    \toprule
    Variables & Coefficient & Std. error & t-value & p-value\\
    \midrule
    \multicolumn{5}{c}{Participant average speed [km/hour]}\\
    \addlinespace[2pt]
    Intercept                                        & 16.023 & 1.621 &  9.883 & \pvl\ ***\\
    Vehicle speed high                               & -7.468 & 0.757 & -9.862 & \pvl\ ***\\
    Pedestrian flow high                             & -3.986 & 0.757 & -5.263 & \pvl\ ***\\
    Bicycle flow high                                & -0.754 & 0.757 & -0.996 & 0.320\\
    Vehicle speed high $\times$ Pedestrian flow high &  4.547 & 0.874 &  5.200 & \pvl\ ***\\
    Vehicle speed high $\times$ Bicycle flow high    &  1.579 & 0.874 &  1.806 & 0.072\\
    Pedestrian flow high $\times$ Bicycle flow high  &  1.000 & 0.874 &  1.144 & 0.254\\
    \midrule
    \multicolumn{2}{@{}l}{Adjusted R-squared} & \multicolumn{3}{l}{0.395}\\
    \bottomrule
  \end{tabular}
  \begin{tablenotes}[flushleft]\footnotesize
    \item \sigstars
  \end{tablenotes}
  \end{threeparttable}
\end{table}

\begin{table}[htbp]
  \centering
  \caption{Model estimation results for perceived unsafety at the unsignalized
    intersection.}
  \label{tab:c13}
  \small
  \begin{threeparttable}
  \begin{tabular}{@{}l r r r l@{}}
    \toprule
    Variables & Odds ratio & Std. error & t-value & p-value\\
    \midrule
    \multicolumn{2}{@{}l}{Intercepts (Thresholds)}
      & \multicolumn{3}{l}{Perceived unsafety [Likert]}\\
    \addlinespace[2pt]
    Perceived unsafety 1 vs 2 (or higher)             & 0.072 & 0.887 & -6.183 & \pvl\ ***\\
    Perceived unsafety 2 (or lower) vs 3 (or higher)  & 0.072 & 0.828 & -6.180 & \pvl\ ***\\
    Perceived unsafety 3 (or lower) vs 4 (or higher)  & 0.540 & 0.818 & -0.354 & 0.723\\
    Perceived unsafety 4 (or lower) vs 5              & 4.059 & 0.822 &  1.704 & 0.088\\
    \addlinespace[2pt]
    Participant average speed high                    & 0.795 & 0.033 & -7.028 & \pvl\ ***\\
    Vehicle speed high                                & 5.468 & 0.438 &  7.883 & \pvl\ ***\\
    Bicycle flow high                                 & 1.377 & 0.374 &  0.999 & 0.317\\
    Pedestrian flow high                              & 0.823 & 0.401 & -0.389 & 0.757\\
    Fixation on pedestrians (\%)                      & 1.020 & 0.014 &  1.369 & 0.171\\
    Fixation on bicycles (\%)                         & 0.868 & 0.044 & -1.495 & 0.135\\
    Fixation on vehicles (\%)                         & 1.055 & 0.020 &  2.716 & 0.007 **\\
    Vehicle speed high $\times$ Pedestrian flow high  & 0.600 & 0.458 & -2.111 & 0.035 *\\
    Vehicle speed high $\times$ Bicycle flow high     & 0.898 & 0.459 & -0.223 & 0.823\\
    Bicycle flow high $\times$ Pedestrian flow high   & 1.843 & 0.437 &  1.403 & 0.162\\
    \midrule
    \multicolumn{2}{@{}l}{Akaike information criterion (AIC)}
      & \multicolumn{3}{l}{834.31}\\
    \bottomrule
  \end{tabular}
  \begin{tablenotes}[flushleft]\footnotesize
    \item \sigstars
  \end{tablenotes}
  \end{threeparttable}
\end{table}

\begin{table}[htbp]
  \centering
  \caption{Model estimation results for fixation on vehicles at the unsignalized
    intersection.}
  \label{tab:c14}
  \small
  \begin{threeparttable}
  \begin{tabular}{@{}l r r r l@{}}
    \toprule
    Variables & Coefficient & Std. error & t-value & p-value\\
    \midrule
    \multicolumn{5}{c}{Fixation on vehicles [\%]}\\
    \addlinespace[2pt]
    Intercept                                        & 10.548 & 2.945 &  3.582 & \pvl\ ***\\
    Participant average speed high                   & -0.238 & 0.087 & -2.730 & 0.007 **\\
    Perceived unsafety high                          & -0.556 & 0.546 & -1.019 & 0.309\\
    Vehicle speed high                               &  4.257 & 1.382 &  3.080 & 0.002 **\\
    Pedestrian flow high                             & -1.489 & 1.282 & -1.162 & 0.246\\
    Bicycle flow high                                & -0.363 & 1.215 & -0.299 & 0.765\\
    Vehicle speed high $\times$ Pedestrian flow high & -0.701 & 1.457 & -0.481 & 0.631\\
    Vehicle speed high $\times$ Bicycle flow high    & -0.175 & 1.424 & -0.123 & 0.902\\
    Pedestrian flow high $\times$ Bicycle flow high  & -0.465 & 1.419 & -0.328 & 0.743\\
    \midrule
    \multicolumn{2}{@{}l}{Adjusted R-squared} & \multicolumn{3}{l}{0.579}\\
    \bottomrule
  \end{tabular}
  \begin{tablenotes}[flushleft]\footnotesize
    \item \sigstars
  \end{tablenotes}
  \end{threeparttable}
\end{table}
